\section{Delivering a control to the bank's Java system}
\label{sec:java}

Java is the required deployment output. The authoring workbench may use Python
for source processing, review and reference evaluation; a transaction must be
able to call the released Java library without running Python or contacting a
language model. The initial integration assumption is a plain Java 17 JAR.
This keeps the decision code independent of Spring, application servers and
particular rule engines. A bank on another Java version needs a declared target
and the same conformance checks before receiving a compatible build.

\subsection{What has actually been built for this proposal}

The delivered \texttt{SPI-Demo1} example contains a source-anchored RuleIR bundle,
synthetic fact snapshots, a separately written Python evaluator, a deterministic
Java emitter, the generated class, two small Java support classes, and a separate
host caller. Its implemented expression subset is deliberately explicit:
Boolean, integer and HK-cent literals and facts; unconditional rule references;
\texttt{all}, \texttt{any}, \texttt{not}; and exact numeric
\texttt{eq}/\texttt{ge}/\texttt{gt} comparisons. The root has a Boolean scope.
Dates, arithmetic expressions, general defaults, conditionals and the full
obligation language remain part of the proposed RuleIR implementation; the demo
compiler rejects them. The demonstration response is its own smaller API profile,
not a claim to implement every field and trace node of the full RuleIR service.

The distinction matters for an implementing agent. The example demonstrates an
actual compiler and integration path that can be extended. It does not silently
turn the 35 general RuleIR contract fixtures into executed tests. Those fixtures
remain the completion target for the full interpreter and backend. The separate
SPI example has 32 executed decision scenarios and eleven consent histories, with
its own recorded evidence.

\subsection{The input and output boundary}

The public API accepts an immutable \texttt{Snapshot}, a time being assessed and
a knowledge cutoff. A \texttt{Snapshot} binds a subject identifier to the
declared fact map. A \texttt{Fact} carries a declared scalar type and one of
\texttt{known}, \texttt{unknown} or \texttt{conflict}. A known Boolean is a Java
\texttt{Boolean}; a known monetary or count value is a \texttt{BigInteger}.
Its evidence identifiers, validity interval and recording time are mandatory.
Unknown has a reason and no value. Conflict identifies competing evidence.

The host supplies every declared field, using an explicit unknown where evidence
is unavailable. Omitting a field or supplying the wrong type is a malformed
request, not a legal conclusion. Validity intervals are half-open: evidence
valid until 10:00 does not establish a fact at 10:00. Evidence recorded after
the knowledge cutoff is unavailable to that replay. The API takes canonical UTC
timestamps with six fractional digits; conversion from a Hong Kong business
timestamp belongs in the host adapter and must preserve the intended instant.

\begin{lstlisting}[language=Java,caption={The library API and one dated monetary input.}]
DecisionRuntime.Result evaluate(
    DecisionRuntime.Snapshot snapshot,
    String validAt,
    String knownAt);

// A known value is exact and supported by dated evidence.
Fact portfolio = Fact.known(
    "money_hkd", new BigInteger("4000000000"),
    List.of("valuation.lee", "ownership.lee"),
    "2026-09-01T00:00:00.000000Z",
    "2026-10-01T00:00:00.000000Z",
    "2026-09-01T00:00:00.000000Z");

Fact missingPortfolio = Fact.unknown("money_hkd", "MISSING");
\end{lstlisting}
The first amount is HK\$40 million expressed in cents. A bank adapter reading a
decimal amount should specify the currency, use exact decimal arithmetic and
reject an unsupported fraction of a cent after the adopted conversion. In Java,
\texttt{amount.movePointRight(2).toBigIntegerExact()} can enforce exact cents
after that policy has been applied. Calling \texttt{doubleValue()} first would
discard the property the specification is trying to preserve. For a currency
conversion, the exchange-rate source, date, direction and rounding rule are part
of the evidence mapping.

The result contains a logical status and an operational disposition:
\begin{center}\small
\begin{tabularx}{\textwidth}{P{0.18\textwidth}P{0.34\textwidth}Y}
\toprule Status & Disposition & Host meaning \\\midrule
\texttt{TRUE} & Conditions met & Selected streamlining conditions established; continue the other controls.\\
\texttt{FALSE} & Standard process required & At least one required condition is refuted; do not apply this streamlined route.\\
\texttt{UNKNOWN} & Review required & Evidence can change the conclusion; show blocking inputs.\\
\texttt{CONFLICT} & Review required & Resolve inconsistent evidence; preserve both records.\\
\texttt{OUT\_OF\_SCOPE} & Outside profile & Select another implemented profile or the standard process.\\
\bottomrule
\end{tabularx}
\end{center}
The demo's JSON includes \texttt{profile}, \texttt{authority}, both timestamps,
the bundle and snapshot hashes, a postorder rule trace, missing and blocking
input lists, and the result hash. A rule-trace entry gives its identifier,
status, source-anchor identifiers and blocking inputs. Source anchors resolve
through the bundle included with the release. Known values and their evidence
remain in the preserved snapshot. The full RuleIR implementation additionally
provides expression-node traces, standardized diagnostics and engine identity
as specified in Section~\ref{sec:product}.

A malformed request throws an explicit Java input exception in this small
library. The production adapter must map that exception, a timeout or an
unavailable release to a technical-failure result and a review/standard-process
path. It must never catch an exception and substitute a successful Boolean.
The demonstration result always carries \texttt{DEMONSTRATION\_ONLY}; production
authority requires a different, approved release record and enforcement by the
host deployment process.

\subsection{How the compiler emits final Java rather than a suggestion}

The emitter first validates the JSON schema, references, identifiers, types and
acyclic rule graph. It then assigns an internal Java method name to each rule
and recursively emits its expression. Identifiers are validated data, never
pieces of executable text supplied by a language model. The generated monetary
comparison is represented by actual Java calls such as:
\begin{lstlisting}[language=Java,caption={The generated financial method's essential body. The complete class is built and tested in the supplied example.}]
return c.rule("spi.financial", sourceIds, () ->
    any(
        compare("ge", c.fact("portfolio"),
            Value.of(new BigInteger("4000000000"))),
        compare("ge", c.fact("net_assets_ex_home"),
            Value.of(new BigInteger("8000000000")))
    ));
\end{lstlisting}
\texttt{compare} returns a known Boolean when both operands are known, or an
unknown with its blocking input names. \texttt{any} implements the truth table
in Equation~\eqref{eq:or}. Java evaluates its arguments before calling the
method, so both branches have the specified eager evaluation and evidence
collection. Replacing this with Java's short-circuit \texttt{||} would change
the trace and missing-input behavior even where the final Boolean agreed.
The context memoizes referenced rules and records their results in a stable
order. A precheck detects conflicting facts in this profile's dependency set
before normal expression evaluation.

The production emitter should implement one mapping per supported operator,
with an explicit error for every unsupported construct. It must not ask an LLM
to write new Java on each regulatory change. The LLM can propose the typed
specification; after review, a deterministic compiler performs the mechanical
translation. This concentrates semantic review on a small language and makes
the generation step reproducible.

For the money comparison, the preservation argument is short. The IR represents
both amounts as mathematical integers of HK cents. The Java boundary constructs
the same integers as \texttt{BigInteger} values. The emitted comparison returns
true exactly when \texttt{compareTo} is nonnegative, which expresses the same
inclusive ordering. The argument depends on the input mapping preserving those
integers. It would fail if the adapter truncated decimals, mixed currencies or
overflowed a narrower integer before constructing the value. For a complete
compiler, analogous local arguments must be composed for every operator,
unknown/error case, rule reference and trace requirement. The delivered tests
check examples of preservation; they are not that universal proof.

\subsection{Reproduce the build and see the host result}

The repository retains a project-local Temurin JDK 17 for this check, verified
against the official archive checksum. It does not change the machine's default
Java installation. From the repository root, the single build command is:

\begin{minipage}{\textwidth}
\begin{lstlisting}
python3 examples/java-dry-run/build_and_verify.py \
  --jdk .localresources/java-toolchain/jdk-17.0.20.1+1
\end{lstlisting}
\end{minipage}

The Python authoring check uses the already available JSON Schema package. The
generated Java library itself has no third-party runtime dependencies. The build
uses \texttt{javac --release 17 -Xlint:all -Werror}; creates the JAR; compiles
the separate fixture and host callers against that JAR; compares complete
reference and Java results; and runs the three mutations. The package has fixed
ZIP timestamps to make repeated builds with the same toolchain reproducible.
The verification record stores compiler identity, source/compiler/JAR digests,
the executed command, results and limitations.

The delivered files below are operational copies of the design, not additional
reading needed to understand why it works:
\begin{center}\small
\begin{tabularx}{\textwidth}{P{0.45\textwidth}Y}
\toprule File under \path{examples/java-dry-run/} & Purpose \\\midrule
\path{spec/spi-control.bundle.json} & Source-linked typed interpretation for the chosen profile.\\
\path{spec/decision-cases.json}; \path{spec/consent-cases.json} & Independent expected statuses and synthetic evidence/history inputs.\\
\path{reference.py}; \path{compile_java.py} & Separate reference evaluation and deterministic Java generation.\\
\path{generated/GeneratedSpi.java} & Final generated decision class.\\
\path{src/main/java/hk/legalmath/spi/} & Evidence, result, three-valued arithmetic/logic support and consent replay.\\
\path{build/spi-controls-demo.jar} & Compiled library callable by a Java host.\\
\path{generated/BankHostExample.java} & Complete synthetic host caller using the public API.\\
\path{build/verification.json}; \path{build/host-output.txt} & Executed validation and host results.\\
\bottomrule
\end{tabularx}
\end{center}

The separate host prints these dispositions for the synthetic history:
\begin{lstlisting}
BEFORE_WITHDRAWAL STREAMLINING_CONDITIONS_MET
AFTER_WITHDRAWAL STANDARD_PROCESS_REQUIRED
RESULT_AUTHORITY DEMONSTRATION_ONLY
\end{lstlisting}
This is the concrete answer to how final code reaches Java: the bank depends on
the released library, constructs the declared evidence snapshot through its
adapter, calls the generated class and handles its structured result. The
compiler, public circulars and drafting tools are build-time/review-time
components. They need not be deployed into the bank's order-processing process.

\subsection{The bank adapter and the exposure race}

The adapter has four responsibilities. It selects the approved rule release for
the legal entity, business profile and relevant time; resolves identities,
categories and evidence into the declared inputs; obtains a consistent consent
and exposure snapshot; and records the result alongside the transaction's other
controls. The generated library should have no credentials for editing these
source systems. A host record should include transaction identifier, release and
snapshot hashes, category and consent generation, data-version identifiers,
decision timestamps and the final disposition.

Concurrency is a material part of correctness. Suppose existing gross exposure
is HK\$8 million and two orders each add HK\$1.5 million against a HK\$10 million
limit. If each reads HK\$8 million, both independently pass at HK\$9.5 million,
but together they produce HK\$11 million. A perfectly correct generated formula
does not prevent that race. The host must include reserved exposure and commit
the new reservation under a transaction or optimistic version check. If the
exposure or consent version changed after evaluation, it reloads and evaluates
again. A withdrawal must invalidate a stale consent version before a new order
can consume the decision. Market-exposure policy, release timing and retry limits
must be agreed with the bank's order and risk owners.

An implementable host sequence is:
\begin{enumerate}
\item Resolve an idempotency key from the order request. A retry with the same
key and unchanged payload returns its recorded disposition; a changed payload
under that key conflicts.
\item Read the approved release and the subject/category consent generation,
holdings, leverage and outstanding reservations with their version identifiers.
Build the dated snapshot and calculate exposure including this order.
\item Evaluate the generated Java library and the host's remaining controls.
If this profile is unresolved or refuted, do not consume its streamlined route.
\item Atomically validate the read versions, reserve the permitted exposure and
append the decision record. On version conflict, repeat from the read step with
a bounded retry policy; otherwise route to review.
\item Retain the committed record for replay. A later correction creates another
assessment linked to the old one; it does not overwrite what was known earlier.
\end{enumerate}
The standalone JAR demonstrates the decision step. It does not simulate this
bank transaction protocol or certify concurrency in an unspecified host system.
Those integration tests are required before production use.

\subsection{Release contents, extension tasks and acceptance}

A production change package must contain the approved source and dependency
register, provision dispositions, rule bundle, reviewer records, supported
profile and effective interval, generated Java source, runtime library, binary
JAR, dependency/toolchain manifest, data-mapping contract, tests and verification
report. The report names the exact JAR tested. The bank's deployment process
verifies these digests and authorization, performs its security and change checks,
then switches the active release atomically. A signature can authenticate the
release authority; a digest alone only identifies content.

The implementation backlog adds a mandatory Java package, W03J, immediately
after the full reference evaluator. It owns the Java domain/API classes, one
emitter per RuleIR operator, result serialization, exact numeric/date behavior,
compiled conformance harness, JAR manifest and host adapter contract. Its pass
condition is equivalence of the complete supported responses on the full
conformance suite, source-derived scenarios and distinguishing mutations,
plus reproducible generation and a separately compiled caller. A backend that
accepts an operator but approximates its meaning fails the package.

W04 must bind release evidence to both the rule bundle and the resulting JAR.
W05 must extend the demonstration's consent replay to the full event contract,
including obligation instances and the adopted calendar policy. W06 must show
the source, interpretation, test case and Java result together. W10 must include
host concurrency, idempotency, stale release and replay tests in its engineering
checks, then separately measure review effort and material errors in the pilot.
This sequencing produces a Java deliverable early while keeping legal
adjudication and production authority explicit.
